% ============================================================================
% ESQUEMA PROPOSTO (banca, sugestão — não mergeado): pré-processamento e os
% DOIS espaços de rótulos da tese (Seção sec:metodo-dados-preproc, 3.2.3).
% A prosa da seção é densa em contagens; a figura mostra a bifurcação que as
% explica: a mesma base normalizada é contada de dois jeitos, com papéis
% distintos — CategorySchema (oráculos) e visão populacional (Cap. 5).
%
% Uso no Cap. 3 (dentro de um ambiente figure):
%   % ============================================================================
% ESQUEMA PROPOSTO (banca, sugestão — não mergeado): pré-processamento e os
% DOIS espaços de rótulos da tese (Seção sec:metodo-dados-preproc, 3.2.3).
% A prosa da seção é densa em contagens; a figura mostra a bifurcação que as
% explica: a mesma base normalizada é contada de dois jeitos, com papéis
% distintos — CategorySchema (oráculos) e visão populacional (Cap. 5).
%
% Uso no Cap. 3 (dentro de um ambiente figure):
%   % ============================================================================
% ESQUEMA PROPOSTO (banca, sugestão — não mergeado): pré-processamento e os
% DOIS espaços de rótulos da tese (Seção sec:metodo-dados-preproc, 3.2.3).
% A prosa da seção é densa em contagens; a figura mostra a bifurcação que as
% explica: a mesma base normalizada é contada de dois jeitos, com papéis
% distintos — CategorySchema (oráculos) e visão populacional (Cap. 5).
%
% Uso no Cap. 3 (dentro de um ambiente figure):
%   % ============================================================================
% ESQUEMA PROPOSTO (banca, sugestão — não mergeado): pré-processamento e os
% DOIS espaços de rótulos da tese (Seção sec:metodo-dados-preproc, 3.2.3).
% A prosa da seção é densa em contagens; a figura mostra a bifurcação que as
% explica: a mesma base normalizada é contada de dois jeitos, com papéis
% distintos — CategorySchema (oráculos) e visão populacional (Cap. 5).
%
% Uso no Cap. 3 (dentro de um ambiente figure):
%   \input{3-metodo/esquemas-propostos/esq-preproc-espacos-rotulos}
% Uso standalone (pré-visualização):
%   \documentclass[tikz,border=6pt]{standalone}
%   \usetikzlibrary{arrows.meta, positioning}
%   \begin{document}\input{esq-preproc-espacos-rotulos.tex}\end{document}
%   (no modo standalone, troque as \ref{...} por texto fixo ou defina-as)
% Requer apenas arrows.meta e positioning (já em packages.tex). Legível em
% P&B: entrada = cinza claro; operação = branco; produto = cinza médio.
% FREEZE respeitado: todos os números são os da própria Seção 3.2.3
% (>=5, >=2, 0,13%, 621, 620, 715, 714, 231.490), nenhum novo.
% ============================================================================
\begin{tikzpicture}[
    x=1cm, y=1cm,
    caixa/.style={draw, rounded corners=2pt, align=center, text width=42mm,
                  minimum height=11mm, inner sep=3pt, font=\footnotesize},
    entrada/.style={caixa, fill=gray!8},
    produto/.style={caixa, fill=gray!14},
    seta/.style={-{Stealth[length=2.4mm]}, thick},
    rot/.style={font=\scriptsize, align=center, inner sep=2pt}
]
    \node[entrada, text width=40mm] (base) at (4.6,0)
        {base corrigida pela auditoria (Seção~\ref{sec:metodo-dados-auditoria}):
         linhas de descrição curta $+$ rótulo};
    \node[caixa, text width=46mm] (norm) at (4.6,-2.4)
        {normalização mínima: minúsculas $+$ remoção de acentos (sem efeito
         nas contagens desta base); rótulos idem; variações lexicais ficam
         com o tokenizador de subpalavras};
    \draw[seta] (base) -- (norm);

    % ------- bifurcação: a mesma base, contada de dois jeitos -------------
    % (o título de cada coluna é a 1ª linha da própria caixa: títulos soltos
    % eram atravessados pelas setas da bifurcação — achado visual)
    \node[caixa, text width=44mm] (filtro5) at (1.4,-5.3)
        {\textbf{espaço dos oráculos} (sobre as linhas): classes com $<5$
         instâncias agrupadas na sentinela \texttt{\_rare\_} ($0{,}13\%$ das
         linhas)};
    \node[caixa, text width=44mm] (dedup) at (8.2,-5.3)
        {\textbf{visão populacional}:\\deduplicação (chave: descrição
         aparada, em minúsculas; mantém a 1ª ocorrência) $+$ filtro brando
         $\ge 2$: 715 classes em linhas, 714 permanecem};

    \node[produto, text width=44mm] (catsch) at (1.4,-8.3)
        {\textit{CategorySchema}: espaço fechado\\
         com \textbf{621 categorias}\\
         (620 mais a sentinela \texttt{\_rare\_});\\
         fonte única da restrição de saída\\
         dos oráculos (Seção~\ref{sec:metodo-oraculo-instrumentacao})};
    \node[produto, text width=44mm] (pop) at (8.2,-8.3)
        {\textbf{231.490 textos únicos} em \textbf{714 classes presentes},
         sem \texttt{\_rare\_} (cauda longa preservada): experimentos do
         Capítulo~\ref{ch:resultados-falco}};

    % setas da bifurcação desenhadas depois dos nós (referência-adiante)
    \draw[seta] (norm.south west) -- (filtro5.north);
    \draw[seta] (norm.south east) -- (dedup.north);
    \draw[seta] (filtro5) -- (catsch);
    \draw[seta] (dedup) -- (pop);

    \node[rot, anchor=west, text width=118mm, align=left] at (-0.9,-10.4)
        {Duas chaves de texto convivem: o censo da auditoria compara
         descrições com os espaços internos colapsados; a deduplicação usa a
         descrição apenas aparada e em minúsculas. A receita executável
         \texttt{scripts/check\_dataset.py} reproduz todas as contagens.};
\end{tikzpicture}

% Uso standalone (pré-visualização):
%   \documentclass[tikz,border=6pt]{standalone}
%   \usetikzlibrary{arrows.meta, positioning}
%   \begin{document}% ============================================================================
% ESQUEMA PROPOSTO (banca, sugestão — não mergeado): pré-processamento e os
% DOIS espaços de rótulos da tese (Seção sec:metodo-dados-preproc, 3.2.3).
% A prosa da seção é densa em contagens; a figura mostra a bifurcação que as
% explica: a mesma base normalizada é contada de dois jeitos, com papéis
% distintos — CategorySchema (oráculos) e visão populacional (Cap. 5).
%
% Uso no Cap. 3 (dentro de um ambiente figure):
%   \input{3-metodo/esquemas-propostos/esq-preproc-espacos-rotulos}
% Uso standalone (pré-visualização):
%   \documentclass[tikz,border=6pt]{standalone}
%   \usetikzlibrary{arrows.meta, positioning}
%   \begin{document}\input{esq-preproc-espacos-rotulos.tex}\end{document}
%   (no modo standalone, troque as \ref{...} por texto fixo ou defina-as)
% Requer apenas arrows.meta e positioning (já em packages.tex). Legível em
% P&B: entrada = cinza claro; operação = branco; produto = cinza médio.
% FREEZE respeitado: todos os números são os da própria Seção 3.2.3
% (>=5, >=2, 0,13%, 621, 620, 715, 714, 231.490), nenhum novo.
% ============================================================================
\begin{tikzpicture}[
    x=1cm, y=1cm,
    caixa/.style={draw, rounded corners=2pt, align=center, text width=42mm,
                  minimum height=11mm, inner sep=3pt, font=\footnotesize},
    entrada/.style={caixa, fill=gray!8},
    produto/.style={caixa, fill=gray!14},
    seta/.style={-{Stealth[length=2.4mm]}, thick},
    rot/.style={font=\scriptsize, align=center, inner sep=2pt}
]
    \node[entrada, text width=40mm] (base) at (4.6,0)
        {base corrigida pela auditoria (Seção~\ref{sec:metodo-dados-auditoria}):
         linhas de descrição curta $+$ rótulo};
    \node[caixa, text width=46mm] (norm) at (4.6,-2.4)
        {normalização mínima: minúsculas $+$ remoção de acentos (sem efeito
         nas contagens desta base); rótulos idem; variações lexicais ficam
         com o tokenizador de subpalavras};
    \draw[seta] (base) -- (norm);

    % ------- bifurcação: a mesma base, contada de dois jeitos -------------
    % (o título de cada coluna é a 1ª linha da própria caixa: títulos soltos
    % eram atravessados pelas setas da bifurcação — achado visual)
    \node[caixa, text width=44mm] (filtro5) at (1.4,-5.3)
        {\textbf{espaço dos oráculos} (sobre as linhas): classes com $<5$
         instâncias agrupadas na sentinela \texttt{\_rare\_} ($0{,}13\%$ das
         linhas)};
    \node[caixa, text width=44mm] (dedup) at (8.2,-5.3)
        {\textbf{visão populacional}:\\deduplicação (chave: descrição
         aparada, em minúsculas; mantém a 1ª ocorrência) $+$ filtro brando
         $\ge 2$: 715 classes em linhas, 714 permanecem};

    \node[produto, text width=44mm] (catsch) at (1.4,-8.3)
        {\textit{CategorySchema}: espaço fechado\\
         com \textbf{621 categorias}\\
         (620 mais a sentinela \texttt{\_rare\_});\\
         fonte única da restrição de saída\\
         dos oráculos (Seção~\ref{sec:metodo-oraculo-instrumentacao})};
    \node[produto, text width=44mm] (pop) at (8.2,-8.3)
        {\textbf{231.490 textos únicos} em \textbf{714 classes presentes},
         sem \texttt{\_rare\_} (cauda longa preservada): experimentos do
         Capítulo~\ref{ch:resultados-falco}};

    % setas da bifurcação desenhadas depois dos nós (referência-adiante)
    \draw[seta] (norm.south west) -- (filtro5.north);
    \draw[seta] (norm.south east) -- (dedup.north);
    \draw[seta] (filtro5) -- (catsch);
    \draw[seta] (dedup) -- (pop);

    \node[rot, anchor=west, text width=118mm, align=left] at (-0.9,-10.4)
        {Duas chaves de texto convivem: o censo da auditoria compara
         descrições com os espaços internos colapsados; a deduplicação usa a
         descrição apenas aparada e em minúsculas. A receita executável
         \texttt{scripts/check\_dataset.py} reproduz todas as contagens.};
\end{tikzpicture}
\end{document}
%   (no modo standalone, troque as \ref{...} por texto fixo ou defina-as)
% Requer apenas arrows.meta e positioning (já em packages.tex). Legível em
% P&B: entrada = cinza claro; operação = branco; produto = cinza médio.
% FREEZE respeitado: todos os números são os da própria Seção 3.2.3
% (>=5, >=2, 0,13%, 621, 620, 715, 714, 231.490), nenhum novo.
% ============================================================================
\begin{tikzpicture}[
    x=1cm, y=1cm,
    caixa/.style={draw, rounded corners=2pt, align=center, text width=42mm,
                  minimum height=11mm, inner sep=3pt, font=\footnotesize},
    entrada/.style={caixa, fill=gray!8},
    produto/.style={caixa, fill=gray!14},
    seta/.style={-{Stealth[length=2.4mm]}, thick},
    rot/.style={font=\scriptsize, align=center, inner sep=2pt}
]
    \node[entrada, text width=40mm] (base) at (4.6,0)
        {base corrigida pela auditoria (Seção~\ref{sec:metodo-dados-auditoria}):
         linhas de descrição curta $+$ rótulo};
    \node[caixa, text width=46mm] (norm) at (4.6,-2.4)
        {normalização mínima: minúsculas $+$ remoção de acentos (sem efeito
         nas contagens desta base); rótulos idem; variações lexicais ficam
         com o tokenizador de subpalavras};
    \draw[seta] (base) -- (norm);

    % ------- bifurcação: a mesma base, contada de dois jeitos -------------
    % (o título de cada coluna é a 1ª linha da própria caixa: títulos soltos
    % eram atravessados pelas setas da bifurcação — achado visual)
    \node[caixa, text width=44mm] (filtro5) at (1.4,-5.3)
        {\textbf{espaço dos oráculos} (sobre as linhas): classes com $<5$
         instâncias agrupadas na sentinela \texttt{\_rare\_} ($0{,}13\%$ das
         linhas)};
    \node[caixa, text width=44mm] (dedup) at (8.2,-5.3)
        {\textbf{visão populacional}:\\deduplicação (chave: descrição
         aparada, em minúsculas; mantém a 1ª ocorrência) $+$ filtro brando
         $\ge 2$: 715 classes em linhas, 714 permanecem};

    \node[produto, text width=44mm] (catsch) at (1.4,-8.3)
        {\textit{CategorySchema}: espaço fechado\\
         com \textbf{621 categorias}\\
         (620 mais a sentinela \texttt{\_rare\_});\\
         fonte única da restrição de saída\\
         dos oráculos (Seção~\ref{sec:metodo-oraculo-instrumentacao})};
    \node[produto, text width=44mm] (pop) at (8.2,-8.3)
        {\textbf{231.490 textos únicos} em \textbf{714 classes presentes},
         sem \texttt{\_rare\_} (cauda longa preservada): experimentos do
         Capítulo~\ref{ch:resultados-falco}};

    % setas da bifurcação desenhadas depois dos nós (referência-adiante)
    \draw[seta] (norm.south west) -- (filtro5.north);
    \draw[seta] (norm.south east) -- (dedup.north);
    \draw[seta] (filtro5) -- (catsch);
    \draw[seta] (dedup) -- (pop);

    \node[rot, anchor=west, text width=118mm, align=left] at (-0.9,-10.4)
        {Duas chaves de texto convivem: o censo da auditoria compara
         descrições com os espaços internos colapsados; a deduplicação usa a
         descrição apenas aparada e em minúsculas. A receita executável
         \texttt{scripts/check\_dataset.py} reproduz todas as contagens.};
\end{tikzpicture}

% Uso standalone (pré-visualização):
%   \documentclass[tikz,border=6pt]{standalone}
%   \usetikzlibrary{arrows.meta, positioning}
%   \begin{document}% ============================================================================
% ESQUEMA PROPOSTO (banca, sugestão — não mergeado): pré-processamento e os
% DOIS espaços de rótulos da tese (Seção sec:metodo-dados-preproc, 3.2.3).
% A prosa da seção é densa em contagens; a figura mostra a bifurcação que as
% explica: a mesma base normalizada é contada de dois jeitos, com papéis
% distintos — CategorySchema (oráculos) e visão populacional (Cap. 5).
%
% Uso no Cap. 3 (dentro de um ambiente figure):
%   % ============================================================================
% ESQUEMA PROPOSTO (banca, sugestão — não mergeado): pré-processamento e os
% DOIS espaços de rótulos da tese (Seção sec:metodo-dados-preproc, 3.2.3).
% A prosa da seção é densa em contagens; a figura mostra a bifurcação que as
% explica: a mesma base normalizada é contada de dois jeitos, com papéis
% distintos — CategorySchema (oráculos) e visão populacional (Cap. 5).
%
% Uso no Cap. 3 (dentro de um ambiente figure):
%   \input{3-metodo/esquemas-propostos/esq-preproc-espacos-rotulos}
% Uso standalone (pré-visualização):
%   \documentclass[tikz,border=6pt]{standalone}
%   \usetikzlibrary{arrows.meta, positioning}
%   \begin{document}\input{esq-preproc-espacos-rotulos.tex}\end{document}
%   (no modo standalone, troque as \ref{...} por texto fixo ou defina-as)
% Requer apenas arrows.meta e positioning (já em packages.tex). Legível em
% P&B: entrada = cinza claro; operação = branco; produto = cinza médio.
% FREEZE respeitado: todos os números são os da própria Seção 3.2.3
% (>=5, >=2, 0,13%, 621, 620, 715, 714, 231.490), nenhum novo.
% ============================================================================
\begin{tikzpicture}[
    x=1cm, y=1cm,
    caixa/.style={draw, rounded corners=2pt, align=center, text width=42mm,
                  minimum height=11mm, inner sep=3pt, font=\footnotesize},
    entrada/.style={caixa, fill=gray!8},
    produto/.style={caixa, fill=gray!14},
    seta/.style={-{Stealth[length=2.4mm]}, thick},
    rot/.style={font=\scriptsize, align=center, inner sep=2pt}
]
    \node[entrada, text width=40mm] (base) at (4.6,0)
        {base corrigida pela auditoria (Seção~\ref{sec:metodo-dados-auditoria}):
         linhas de descrição curta $+$ rótulo};
    \node[caixa, text width=46mm] (norm) at (4.6,-2.4)
        {normalização mínima: minúsculas $+$ remoção de acentos (sem efeito
         nas contagens desta base); rótulos idem; variações lexicais ficam
         com o tokenizador de subpalavras};
    \draw[seta] (base) -- (norm);

    % ------- bifurcação: a mesma base, contada de dois jeitos -------------
    % (o título de cada coluna é a 1ª linha da própria caixa: títulos soltos
    % eram atravessados pelas setas da bifurcação — achado visual)
    \node[caixa, text width=44mm] (filtro5) at (1.4,-5.3)
        {\textbf{espaço dos oráculos} (sobre as linhas): classes com $<5$
         instâncias agrupadas na sentinela \texttt{\_rare\_} ($0{,}13\%$ das
         linhas)};
    \node[caixa, text width=44mm] (dedup) at (8.2,-5.3)
        {\textbf{visão populacional}:\\deduplicação (chave: descrição
         aparada, em minúsculas; mantém a 1ª ocorrência) $+$ filtro brando
         $\ge 2$: 715 classes em linhas, 714 permanecem};

    \node[produto, text width=44mm] (catsch) at (1.4,-8.3)
        {\textit{CategorySchema}: espaço fechado\\
         com \textbf{621 categorias}\\
         (620 mais a sentinela \texttt{\_rare\_});\\
         fonte única da restrição de saída\\
         dos oráculos (Seção~\ref{sec:metodo-oraculo-instrumentacao})};
    \node[produto, text width=44mm] (pop) at (8.2,-8.3)
        {\textbf{231.490 textos únicos} em \textbf{714 classes presentes},
         sem \texttt{\_rare\_} (cauda longa preservada): experimentos do
         Capítulo~\ref{ch:resultados-falco}};

    % setas da bifurcação desenhadas depois dos nós (referência-adiante)
    \draw[seta] (norm.south west) -- (filtro5.north);
    \draw[seta] (norm.south east) -- (dedup.north);
    \draw[seta] (filtro5) -- (catsch);
    \draw[seta] (dedup) -- (pop);

    \node[rot, anchor=west, text width=118mm, align=left] at (-0.9,-10.4)
        {Duas chaves de texto convivem: o censo da auditoria compara
         descrições com os espaços internos colapsados; a deduplicação usa a
         descrição apenas aparada e em minúsculas. A receita executável
         \texttt{scripts/check\_dataset.py} reproduz todas as contagens.};
\end{tikzpicture}

% Uso standalone (pré-visualização):
%   \documentclass[tikz,border=6pt]{standalone}
%   \usetikzlibrary{arrows.meta, positioning}
%   \begin{document}% ============================================================================
% ESQUEMA PROPOSTO (banca, sugestão — não mergeado): pré-processamento e os
% DOIS espaços de rótulos da tese (Seção sec:metodo-dados-preproc, 3.2.3).
% A prosa da seção é densa em contagens; a figura mostra a bifurcação que as
% explica: a mesma base normalizada é contada de dois jeitos, com papéis
% distintos — CategorySchema (oráculos) e visão populacional (Cap. 5).
%
% Uso no Cap. 3 (dentro de um ambiente figure):
%   \input{3-metodo/esquemas-propostos/esq-preproc-espacos-rotulos}
% Uso standalone (pré-visualização):
%   \documentclass[tikz,border=6pt]{standalone}
%   \usetikzlibrary{arrows.meta, positioning}
%   \begin{document}\input{esq-preproc-espacos-rotulos.tex}\end{document}
%   (no modo standalone, troque as \ref{...} por texto fixo ou defina-as)
% Requer apenas arrows.meta e positioning (já em packages.tex). Legível em
% P&B: entrada = cinza claro; operação = branco; produto = cinza médio.
% FREEZE respeitado: todos os números são os da própria Seção 3.2.3
% (>=5, >=2, 0,13%, 621, 620, 715, 714, 231.490), nenhum novo.
% ============================================================================
\begin{tikzpicture}[
    x=1cm, y=1cm,
    caixa/.style={draw, rounded corners=2pt, align=center, text width=42mm,
                  minimum height=11mm, inner sep=3pt, font=\footnotesize},
    entrada/.style={caixa, fill=gray!8},
    produto/.style={caixa, fill=gray!14},
    seta/.style={-{Stealth[length=2.4mm]}, thick},
    rot/.style={font=\scriptsize, align=center, inner sep=2pt}
]
    \node[entrada, text width=40mm] (base) at (4.6,0)
        {base corrigida pela auditoria (Seção~\ref{sec:metodo-dados-auditoria}):
         linhas de descrição curta $+$ rótulo};
    \node[caixa, text width=46mm] (norm) at (4.6,-2.4)
        {normalização mínima: minúsculas $+$ remoção de acentos (sem efeito
         nas contagens desta base); rótulos idem; variações lexicais ficam
         com o tokenizador de subpalavras};
    \draw[seta] (base) -- (norm);

    % ------- bifurcação: a mesma base, contada de dois jeitos -------------
    % (o título de cada coluna é a 1ª linha da própria caixa: títulos soltos
    % eram atravessados pelas setas da bifurcação — achado visual)
    \node[caixa, text width=44mm] (filtro5) at (1.4,-5.3)
        {\textbf{espaço dos oráculos} (sobre as linhas): classes com $<5$
         instâncias agrupadas na sentinela \texttt{\_rare\_} ($0{,}13\%$ das
         linhas)};
    \node[caixa, text width=44mm] (dedup) at (8.2,-5.3)
        {\textbf{visão populacional}:\\deduplicação (chave: descrição
         aparada, em minúsculas; mantém a 1ª ocorrência) $+$ filtro brando
         $\ge 2$: 715 classes em linhas, 714 permanecem};

    \node[produto, text width=44mm] (catsch) at (1.4,-8.3)
        {\textit{CategorySchema}: espaço fechado\\
         com \textbf{621 categorias}\\
         (620 mais a sentinela \texttt{\_rare\_});\\
         fonte única da restrição de saída\\
         dos oráculos (Seção~\ref{sec:metodo-oraculo-instrumentacao})};
    \node[produto, text width=44mm] (pop) at (8.2,-8.3)
        {\textbf{231.490 textos únicos} em \textbf{714 classes presentes},
         sem \texttt{\_rare\_} (cauda longa preservada): experimentos do
         Capítulo~\ref{ch:resultados-falco}};

    % setas da bifurcação desenhadas depois dos nós (referência-adiante)
    \draw[seta] (norm.south west) -- (filtro5.north);
    \draw[seta] (norm.south east) -- (dedup.north);
    \draw[seta] (filtro5) -- (catsch);
    \draw[seta] (dedup) -- (pop);

    \node[rot, anchor=west, text width=118mm, align=left] at (-0.9,-10.4)
        {Duas chaves de texto convivem: o censo da auditoria compara
         descrições com os espaços internos colapsados; a deduplicação usa a
         descrição apenas aparada e em minúsculas. A receita executável
         \texttt{scripts/check\_dataset.py} reproduz todas as contagens.};
\end{tikzpicture}
\end{document}
%   (no modo standalone, troque as \ref{...} por texto fixo ou defina-as)
% Requer apenas arrows.meta e positioning (já em packages.tex). Legível em
% P&B: entrada = cinza claro; operação = branco; produto = cinza médio.
% FREEZE respeitado: todos os números são os da própria Seção 3.2.3
% (>=5, >=2, 0,13%, 621, 620, 715, 714, 231.490), nenhum novo.
% ============================================================================
\begin{tikzpicture}[
    x=1cm, y=1cm,
    caixa/.style={draw, rounded corners=2pt, align=center, text width=42mm,
                  minimum height=11mm, inner sep=3pt, font=\footnotesize},
    entrada/.style={caixa, fill=gray!8},
    produto/.style={caixa, fill=gray!14},
    seta/.style={-{Stealth[length=2.4mm]}, thick},
    rot/.style={font=\scriptsize, align=center, inner sep=2pt}
]
    \node[entrada, text width=40mm] (base) at (4.6,0)
        {base corrigida pela auditoria (Seção~\ref{sec:metodo-dados-auditoria}):
         linhas de descrição curta $+$ rótulo};
    \node[caixa, text width=46mm] (norm) at (4.6,-2.4)
        {normalização mínima: minúsculas $+$ remoção de acentos (sem efeito
         nas contagens desta base); rótulos idem; variações lexicais ficam
         com o tokenizador de subpalavras};
    \draw[seta] (base) -- (norm);

    % ------- bifurcação: a mesma base, contada de dois jeitos -------------
    % (o título de cada coluna é a 1ª linha da própria caixa: títulos soltos
    % eram atravessados pelas setas da bifurcação — achado visual)
    \node[caixa, text width=44mm] (filtro5) at (1.4,-5.3)
        {\textbf{espaço dos oráculos} (sobre as linhas): classes com $<5$
         instâncias agrupadas na sentinela \texttt{\_rare\_} ($0{,}13\%$ das
         linhas)};
    \node[caixa, text width=44mm] (dedup) at (8.2,-5.3)
        {\textbf{visão populacional}:\\deduplicação (chave: descrição
         aparada, em minúsculas; mantém a 1ª ocorrência) $+$ filtro brando
         $\ge 2$: 715 classes em linhas, 714 permanecem};

    \node[produto, text width=44mm] (catsch) at (1.4,-8.3)
        {\textit{CategorySchema}: espaço fechado\\
         com \textbf{621 categorias}\\
         (620 mais a sentinela \texttt{\_rare\_});\\
         fonte única da restrição de saída\\
         dos oráculos (Seção~\ref{sec:metodo-oraculo-instrumentacao})};
    \node[produto, text width=44mm] (pop) at (8.2,-8.3)
        {\textbf{231.490 textos únicos} em \textbf{714 classes presentes},
         sem \texttt{\_rare\_} (cauda longa preservada): experimentos do
         Capítulo~\ref{ch:resultados-falco}};

    % setas da bifurcação desenhadas depois dos nós (referência-adiante)
    \draw[seta] (norm.south west) -- (filtro5.north);
    \draw[seta] (norm.south east) -- (dedup.north);
    \draw[seta] (filtro5) -- (catsch);
    \draw[seta] (dedup) -- (pop);

    \node[rot, anchor=west, text width=118mm, align=left] at (-0.9,-10.4)
        {Duas chaves de texto convivem: o censo da auditoria compara
         descrições com os espaços internos colapsados; a deduplicação usa a
         descrição apenas aparada e em minúsculas. A receita executável
         \texttt{scripts/check\_dataset.py} reproduz todas as contagens.};
\end{tikzpicture}
\end{document}
%   (no modo standalone, troque as \ref{...} por texto fixo ou defina-as)
% Requer apenas arrows.meta e positioning (já em packages.tex). Legível em
% P&B: entrada = cinza claro; operação = branco; produto = cinza médio.
% FREEZE respeitado: todos os números são os da própria Seção 3.2.3
% (>=5, >=2, 0,13%, 621, 620, 715, 714, 231.490), nenhum novo.
% ============================================================================
\begin{tikzpicture}[
    x=1cm, y=1cm,
    caixa/.style={draw, rounded corners=2pt, align=center, text width=42mm,
                  minimum height=11mm, inner sep=3pt, font=\footnotesize},
    entrada/.style={caixa, fill=gray!8},
    produto/.style={caixa, fill=gray!14},
    seta/.style={-{Stealth[length=2.4mm]}, thick},
    rot/.style={font=\scriptsize, align=center, inner sep=2pt}
]
    \node[entrada, text width=40mm] (base) at (4.6,0)
        {base corrigida pela auditoria (Seção~\ref{sec:metodo-dados-auditoria}):
         linhas de descrição curta $+$ rótulo};
    \node[caixa, text width=46mm] (norm) at (4.6,-2.4)
        {normalização mínima: minúsculas $+$ remoção de acentos (sem efeito
         nas contagens desta base); rótulos idem; variações lexicais ficam
         com o tokenizador de subpalavras};
    \draw[seta] (base) -- (norm);

    % ------- bifurcação: a mesma base, contada de dois jeitos -------------
    % (o título de cada coluna é a 1ª linha da própria caixa: títulos soltos
    % eram atravessados pelas setas da bifurcação — achado visual)
    \node[caixa, text width=44mm] (filtro5) at (1.4,-5.3)
        {\textbf{espaço dos oráculos} (sobre as linhas): classes com $<5$
         instâncias agrupadas na sentinela \texttt{\_rare\_} ($0{,}13\%$ das
         linhas)};
    \node[caixa, text width=44mm] (dedup) at (8.2,-5.3)
        {\textbf{visão populacional}:\\deduplicação (chave: descrição
         aparada, em minúsculas; mantém a 1ª ocorrência) $+$ filtro brando
         $\ge 2$: 715 classes em linhas, 714 permanecem};

    \node[produto, text width=44mm] (catsch) at (1.4,-8.3)
        {\textit{CategorySchema}: espaço fechado\\
         com \textbf{621 categorias}\\
         (620 mais a sentinela \texttt{\_rare\_});\\
         fonte única da restrição de saída\\
         dos oráculos (Seção~\ref{sec:metodo-oraculo-instrumentacao})};
    \node[produto, text width=44mm] (pop) at (8.2,-8.3)
        {\textbf{231.490 textos únicos} em \textbf{714 classes presentes},
         sem \texttt{\_rare\_} (cauda longa preservada): experimentos do
         Capítulo~\ref{ch:resultados-falco}};

    % setas da bifurcação desenhadas depois dos nós (referência-adiante)
    \draw[seta] (norm.south west) -- (filtro5.north);
    \draw[seta] (norm.south east) -- (dedup.north);
    \draw[seta] (filtro5) -- (catsch);
    \draw[seta] (dedup) -- (pop);

    \node[rot, anchor=west, text width=118mm, align=left] at (-0.9,-10.4)
        {Duas chaves de texto convivem: o censo da auditoria compara
         descrições com os espaços internos colapsados; a deduplicação usa a
         descrição apenas aparada e em minúsculas. A receita executável
         \texttt{scripts/check\_dataset.py} reproduz todas as contagens.};
\end{tikzpicture}

% Uso standalone (pré-visualização):
%   \documentclass[tikz,border=6pt]{standalone}
%   \usetikzlibrary{arrows.meta, positioning}
%   \begin{document}% ============================================================================
% ESQUEMA PROPOSTO (banca, sugestão — não mergeado): pré-processamento e os
% DOIS espaços de rótulos da tese (Seção sec:metodo-dados-preproc, 3.2.3).
% A prosa da seção é densa em contagens; a figura mostra a bifurcação que as
% explica: a mesma base normalizada é contada de dois jeitos, com papéis
% distintos — CategorySchema (oráculos) e visão populacional (Cap. 5).
%
% Uso no Cap. 3 (dentro de um ambiente figure):
%   % ============================================================================
% ESQUEMA PROPOSTO (banca, sugestão — não mergeado): pré-processamento e os
% DOIS espaços de rótulos da tese (Seção sec:metodo-dados-preproc, 3.2.3).
% A prosa da seção é densa em contagens; a figura mostra a bifurcação que as
% explica: a mesma base normalizada é contada de dois jeitos, com papéis
% distintos — CategorySchema (oráculos) e visão populacional (Cap. 5).
%
% Uso no Cap. 3 (dentro de um ambiente figure):
%   % ============================================================================
% ESQUEMA PROPOSTO (banca, sugestão — não mergeado): pré-processamento e os
% DOIS espaços de rótulos da tese (Seção sec:metodo-dados-preproc, 3.2.3).
% A prosa da seção é densa em contagens; a figura mostra a bifurcação que as
% explica: a mesma base normalizada é contada de dois jeitos, com papéis
% distintos — CategorySchema (oráculos) e visão populacional (Cap. 5).
%
% Uso no Cap. 3 (dentro de um ambiente figure):
%   \input{3-metodo/esquemas-propostos/esq-preproc-espacos-rotulos}
% Uso standalone (pré-visualização):
%   \documentclass[tikz,border=6pt]{standalone}
%   \usetikzlibrary{arrows.meta, positioning}
%   \begin{document}\input{esq-preproc-espacos-rotulos.tex}\end{document}
%   (no modo standalone, troque as \ref{...} por texto fixo ou defina-as)
% Requer apenas arrows.meta e positioning (já em packages.tex). Legível em
% P&B: entrada = cinza claro; operação = branco; produto = cinza médio.
% FREEZE respeitado: todos os números são os da própria Seção 3.2.3
% (>=5, >=2, 0,13%, 621, 620, 715, 714, 231.490), nenhum novo.
% ============================================================================
\begin{tikzpicture}[
    x=1cm, y=1cm,
    caixa/.style={draw, rounded corners=2pt, align=center, text width=42mm,
                  minimum height=11mm, inner sep=3pt, font=\footnotesize},
    entrada/.style={caixa, fill=gray!8},
    produto/.style={caixa, fill=gray!14},
    seta/.style={-{Stealth[length=2.4mm]}, thick},
    rot/.style={font=\scriptsize, align=center, inner sep=2pt}
]
    \node[entrada, text width=40mm] (base) at (4.6,0)
        {base corrigida pela auditoria (Seção~\ref{sec:metodo-dados-auditoria}):
         linhas de descrição curta $+$ rótulo};
    \node[caixa, text width=46mm] (norm) at (4.6,-2.4)
        {normalização mínima: minúsculas $+$ remoção de acentos (sem efeito
         nas contagens desta base); rótulos idem; variações lexicais ficam
         com o tokenizador de subpalavras};
    \draw[seta] (base) -- (norm);

    % ------- bifurcação: a mesma base, contada de dois jeitos -------------
    % (o título de cada coluna é a 1ª linha da própria caixa: títulos soltos
    % eram atravessados pelas setas da bifurcação — achado visual)
    \node[caixa, text width=44mm] (filtro5) at (1.4,-5.3)
        {\textbf{espaço dos oráculos} (sobre as linhas): classes com $<5$
         instâncias agrupadas na sentinela \texttt{\_rare\_} ($0{,}13\%$ das
         linhas)};
    \node[caixa, text width=44mm] (dedup) at (8.2,-5.3)
        {\textbf{visão populacional}:\\deduplicação (chave: descrição
         aparada, em minúsculas; mantém a 1ª ocorrência) $+$ filtro brando
         $\ge 2$: 715 classes em linhas, 714 permanecem};

    \node[produto, text width=44mm] (catsch) at (1.4,-8.3)
        {\textit{CategorySchema}: espaço fechado\\
         com \textbf{621 categorias}\\
         (620 mais a sentinela \texttt{\_rare\_});\\
         fonte única da restrição de saída\\
         dos oráculos (Seção~\ref{sec:metodo-oraculo-instrumentacao})};
    \node[produto, text width=44mm] (pop) at (8.2,-8.3)
        {\textbf{231.490 textos únicos} em \textbf{714 classes presentes},
         sem \texttt{\_rare\_} (cauda longa preservada): experimentos do
         Capítulo~\ref{ch:resultados-falco}};

    % setas da bifurcação desenhadas depois dos nós (referência-adiante)
    \draw[seta] (norm.south west) -- (filtro5.north);
    \draw[seta] (norm.south east) -- (dedup.north);
    \draw[seta] (filtro5) -- (catsch);
    \draw[seta] (dedup) -- (pop);

    \node[rot, anchor=west, text width=118mm, align=left] at (-0.9,-10.4)
        {Duas chaves de texto convivem: o censo da auditoria compara
         descrições com os espaços internos colapsados; a deduplicação usa a
         descrição apenas aparada e em minúsculas. A receita executável
         \texttt{scripts/check\_dataset.py} reproduz todas as contagens.};
\end{tikzpicture}

% Uso standalone (pré-visualização):
%   \documentclass[tikz,border=6pt]{standalone}
%   \usetikzlibrary{arrows.meta, positioning}
%   \begin{document}% ============================================================================
% ESQUEMA PROPOSTO (banca, sugestão — não mergeado): pré-processamento e os
% DOIS espaços de rótulos da tese (Seção sec:metodo-dados-preproc, 3.2.3).
% A prosa da seção é densa em contagens; a figura mostra a bifurcação que as
% explica: a mesma base normalizada é contada de dois jeitos, com papéis
% distintos — CategorySchema (oráculos) e visão populacional (Cap. 5).
%
% Uso no Cap. 3 (dentro de um ambiente figure):
%   \input{3-metodo/esquemas-propostos/esq-preproc-espacos-rotulos}
% Uso standalone (pré-visualização):
%   \documentclass[tikz,border=6pt]{standalone}
%   \usetikzlibrary{arrows.meta, positioning}
%   \begin{document}\input{esq-preproc-espacos-rotulos.tex}\end{document}
%   (no modo standalone, troque as \ref{...} por texto fixo ou defina-as)
% Requer apenas arrows.meta e positioning (já em packages.tex). Legível em
% P&B: entrada = cinza claro; operação = branco; produto = cinza médio.
% FREEZE respeitado: todos os números são os da própria Seção 3.2.3
% (>=5, >=2, 0,13%, 621, 620, 715, 714, 231.490), nenhum novo.
% ============================================================================
\begin{tikzpicture}[
    x=1cm, y=1cm,
    caixa/.style={draw, rounded corners=2pt, align=center, text width=42mm,
                  minimum height=11mm, inner sep=3pt, font=\footnotesize},
    entrada/.style={caixa, fill=gray!8},
    produto/.style={caixa, fill=gray!14},
    seta/.style={-{Stealth[length=2.4mm]}, thick},
    rot/.style={font=\scriptsize, align=center, inner sep=2pt}
]
    \node[entrada, text width=40mm] (base) at (4.6,0)
        {base corrigida pela auditoria (Seção~\ref{sec:metodo-dados-auditoria}):
         linhas de descrição curta $+$ rótulo};
    \node[caixa, text width=46mm] (norm) at (4.6,-2.4)
        {normalização mínima: minúsculas $+$ remoção de acentos (sem efeito
         nas contagens desta base); rótulos idem; variações lexicais ficam
         com o tokenizador de subpalavras};
    \draw[seta] (base) -- (norm);

    % ------- bifurcação: a mesma base, contada de dois jeitos -------------
    % (o título de cada coluna é a 1ª linha da própria caixa: títulos soltos
    % eram atravessados pelas setas da bifurcação — achado visual)
    \node[caixa, text width=44mm] (filtro5) at (1.4,-5.3)
        {\textbf{espaço dos oráculos} (sobre as linhas): classes com $<5$
         instâncias agrupadas na sentinela \texttt{\_rare\_} ($0{,}13\%$ das
         linhas)};
    \node[caixa, text width=44mm] (dedup) at (8.2,-5.3)
        {\textbf{visão populacional}:\\deduplicação (chave: descrição
         aparada, em minúsculas; mantém a 1ª ocorrência) $+$ filtro brando
         $\ge 2$: 715 classes em linhas, 714 permanecem};

    \node[produto, text width=44mm] (catsch) at (1.4,-8.3)
        {\textit{CategorySchema}: espaço fechado\\
         com \textbf{621 categorias}\\
         (620 mais a sentinela \texttt{\_rare\_});\\
         fonte única da restrição de saída\\
         dos oráculos (Seção~\ref{sec:metodo-oraculo-instrumentacao})};
    \node[produto, text width=44mm] (pop) at (8.2,-8.3)
        {\textbf{231.490 textos únicos} em \textbf{714 classes presentes},
         sem \texttt{\_rare\_} (cauda longa preservada): experimentos do
         Capítulo~\ref{ch:resultados-falco}};

    % setas da bifurcação desenhadas depois dos nós (referência-adiante)
    \draw[seta] (norm.south west) -- (filtro5.north);
    \draw[seta] (norm.south east) -- (dedup.north);
    \draw[seta] (filtro5) -- (catsch);
    \draw[seta] (dedup) -- (pop);

    \node[rot, anchor=west, text width=118mm, align=left] at (-0.9,-10.4)
        {Duas chaves de texto convivem: o censo da auditoria compara
         descrições com os espaços internos colapsados; a deduplicação usa a
         descrição apenas aparada e em minúsculas. A receita executável
         \texttt{scripts/check\_dataset.py} reproduz todas as contagens.};
\end{tikzpicture}
\end{document}
%   (no modo standalone, troque as \ref{...} por texto fixo ou defina-as)
% Requer apenas arrows.meta e positioning (já em packages.tex). Legível em
% P&B: entrada = cinza claro; operação = branco; produto = cinza médio.
% FREEZE respeitado: todos os números são os da própria Seção 3.2.3
% (>=5, >=2, 0,13%, 621, 620, 715, 714, 231.490), nenhum novo.
% ============================================================================
\begin{tikzpicture}[
    x=1cm, y=1cm,
    caixa/.style={draw, rounded corners=2pt, align=center, text width=42mm,
                  minimum height=11mm, inner sep=3pt, font=\footnotesize},
    entrada/.style={caixa, fill=gray!8},
    produto/.style={caixa, fill=gray!14},
    seta/.style={-{Stealth[length=2.4mm]}, thick},
    rot/.style={font=\scriptsize, align=center, inner sep=2pt}
]
    \node[entrada, text width=40mm] (base) at (4.6,0)
        {base corrigida pela auditoria (Seção~\ref{sec:metodo-dados-auditoria}):
         linhas de descrição curta $+$ rótulo};
    \node[caixa, text width=46mm] (norm) at (4.6,-2.4)
        {normalização mínima: minúsculas $+$ remoção de acentos (sem efeito
         nas contagens desta base); rótulos idem; variações lexicais ficam
         com o tokenizador de subpalavras};
    \draw[seta] (base) -- (norm);

    % ------- bifurcação: a mesma base, contada de dois jeitos -------------
    % (o título de cada coluna é a 1ª linha da própria caixa: títulos soltos
    % eram atravessados pelas setas da bifurcação — achado visual)
    \node[caixa, text width=44mm] (filtro5) at (1.4,-5.3)
        {\textbf{espaço dos oráculos} (sobre as linhas): classes com $<5$
         instâncias agrupadas na sentinela \texttt{\_rare\_} ($0{,}13\%$ das
         linhas)};
    \node[caixa, text width=44mm] (dedup) at (8.2,-5.3)
        {\textbf{visão populacional}:\\deduplicação (chave: descrição
         aparada, em minúsculas; mantém a 1ª ocorrência) $+$ filtro brando
         $\ge 2$: 715 classes em linhas, 714 permanecem};

    \node[produto, text width=44mm] (catsch) at (1.4,-8.3)
        {\textit{CategorySchema}: espaço fechado\\
         com \textbf{621 categorias}\\
         (620 mais a sentinela \texttt{\_rare\_});\\
         fonte única da restrição de saída\\
         dos oráculos (Seção~\ref{sec:metodo-oraculo-instrumentacao})};
    \node[produto, text width=44mm] (pop) at (8.2,-8.3)
        {\textbf{231.490 textos únicos} em \textbf{714 classes presentes},
         sem \texttt{\_rare\_} (cauda longa preservada): experimentos do
         Capítulo~\ref{ch:resultados-falco}};

    % setas da bifurcação desenhadas depois dos nós (referência-adiante)
    \draw[seta] (norm.south west) -- (filtro5.north);
    \draw[seta] (norm.south east) -- (dedup.north);
    \draw[seta] (filtro5) -- (catsch);
    \draw[seta] (dedup) -- (pop);

    \node[rot, anchor=west, text width=118mm, align=left] at (-0.9,-10.4)
        {Duas chaves de texto convivem: o censo da auditoria compara
         descrições com os espaços internos colapsados; a deduplicação usa a
         descrição apenas aparada e em minúsculas. A receita executável
         \texttt{scripts/check\_dataset.py} reproduz todas as contagens.};
\end{tikzpicture}

% Uso standalone (pré-visualização):
%   \documentclass[tikz,border=6pt]{standalone}
%   \usetikzlibrary{arrows.meta, positioning}
%   \begin{document}% ============================================================================
% ESQUEMA PROPOSTO (banca, sugestão — não mergeado): pré-processamento e os
% DOIS espaços de rótulos da tese (Seção sec:metodo-dados-preproc, 3.2.3).
% A prosa da seção é densa em contagens; a figura mostra a bifurcação que as
% explica: a mesma base normalizada é contada de dois jeitos, com papéis
% distintos — CategorySchema (oráculos) e visão populacional (Cap. 5).
%
% Uso no Cap. 3 (dentro de um ambiente figure):
%   % ============================================================================
% ESQUEMA PROPOSTO (banca, sugestão — não mergeado): pré-processamento e os
% DOIS espaços de rótulos da tese (Seção sec:metodo-dados-preproc, 3.2.3).
% A prosa da seção é densa em contagens; a figura mostra a bifurcação que as
% explica: a mesma base normalizada é contada de dois jeitos, com papéis
% distintos — CategorySchema (oráculos) e visão populacional (Cap. 5).
%
% Uso no Cap. 3 (dentro de um ambiente figure):
%   \input{3-metodo/esquemas-propostos/esq-preproc-espacos-rotulos}
% Uso standalone (pré-visualização):
%   \documentclass[tikz,border=6pt]{standalone}
%   \usetikzlibrary{arrows.meta, positioning}
%   \begin{document}\input{esq-preproc-espacos-rotulos.tex}\end{document}
%   (no modo standalone, troque as \ref{...} por texto fixo ou defina-as)
% Requer apenas arrows.meta e positioning (já em packages.tex). Legível em
% P&B: entrada = cinza claro; operação = branco; produto = cinza médio.
% FREEZE respeitado: todos os números são os da própria Seção 3.2.3
% (>=5, >=2, 0,13%, 621, 620, 715, 714, 231.490), nenhum novo.
% ============================================================================
\begin{tikzpicture}[
    x=1cm, y=1cm,
    caixa/.style={draw, rounded corners=2pt, align=center, text width=42mm,
                  minimum height=11mm, inner sep=3pt, font=\footnotesize},
    entrada/.style={caixa, fill=gray!8},
    produto/.style={caixa, fill=gray!14},
    seta/.style={-{Stealth[length=2.4mm]}, thick},
    rot/.style={font=\scriptsize, align=center, inner sep=2pt}
]
    \node[entrada, text width=40mm] (base) at (4.6,0)
        {base corrigida pela auditoria (Seção~\ref{sec:metodo-dados-auditoria}):
         linhas de descrição curta $+$ rótulo};
    \node[caixa, text width=46mm] (norm) at (4.6,-2.4)
        {normalização mínima: minúsculas $+$ remoção de acentos (sem efeito
         nas contagens desta base); rótulos idem; variações lexicais ficam
         com o tokenizador de subpalavras};
    \draw[seta] (base) -- (norm);

    % ------- bifurcação: a mesma base, contada de dois jeitos -------------
    % (o título de cada coluna é a 1ª linha da própria caixa: títulos soltos
    % eram atravessados pelas setas da bifurcação — achado visual)
    \node[caixa, text width=44mm] (filtro5) at (1.4,-5.3)
        {\textbf{espaço dos oráculos} (sobre as linhas): classes com $<5$
         instâncias agrupadas na sentinela \texttt{\_rare\_} ($0{,}13\%$ das
         linhas)};
    \node[caixa, text width=44mm] (dedup) at (8.2,-5.3)
        {\textbf{visão populacional}:\\deduplicação (chave: descrição
         aparada, em minúsculas; mantém a 1ª ocorrência) $+$ filtro brando
         $\ge 2$: 715 classes em linhas, 714 permanecem};

    \node[produto, text width=44mm] (catsch) at (1.4,-8.3)
        {\textit{CategorySchema}: espaço fechado\\
         com \textbf{621 categorias}\\
         (620 mais a sentinela \texttt{\_rare\_});\\
         fonte única da restrição de saída\\
         dos oráculos (Seção~\ref{sec:metodo-oraculo-instrumentacao})};
    \node[produto, text width=44mm] (pop) at (8.2,-8.3)
        {\textbf{231.490 textos únicos} em \textbf{714 classes presentes},
         sem \texttt{\_rare\_} (cauda longa preservada): experimentos do
         Capítulo~\ref{ch:resultados-falco}};

    % setas da bifurcação desenhadas depois dos nós (referência-adiante)
    \draw[seta] (norm.south west) -- (filtro5.north);
    \draw[seta] (norm.south east) -- (dedup.north);
    \draw[seta] (filtro5) -- (catsch);
    \draw[seta] (dedup) -- (pop);

    \node[rot, anchor=west, text width=118mm, align=left] at (-0.9,-10.4)
        {Duas chaves de texto convivem: o censo da auditoria compara
         descrições com os espaços internos colapsados; a deduplicação usa a
         descrição apenas aparada e em minúsculas. A receita executável
         \texttt{scripts/check\_dataset.py} reproduz todas as contagens.};
\end{tikzpicture}

% Uso standalone (pré-visualização):
%   \documentclass[tikz,border=6pt]{standalone}
%   \usetikzlibrary{arrows.meta, positioning}
%   \begin{document}% ============================================================================
% ESQUEMA PROPOSTO (banca, sugestão — não mergeado): pré-processamento e os
% DOIS espaços de rótulos da tese (Seção sec:metodo-dados-preproc, 3.2.3).
% A prosa da seção é densa em contagens; a figura mostra a bifurcação que as
% explica: a mesma base normalizada é contada de dois jeitos, com papéis
% distintos — CategorySchema (oráculos) e visão populacional (Cap. 5).
%
% Uso no Cap. 3 (dentro de um ambiente figure):
%   \input{3-metodo/esquemas-propostos/esq-preproc-espacos-rotulos}
% Uso standalone (pré-visualização):
%   \documentclass[tikz,border=6pt]{standalone}
%   \usetikzlibrary{arrows.meta, positioning}
%   \begin{document}\input{esq-preproc-espacos-rotulos.tex}\end{document}
%   (no modo standalone, troque as \ref{...} por texto fixo ou defina-as)
% Requer apenas arrows.meta e positioning (já em packages.tex). Legível em
% P&B: entrada = cinza claro; operação = branco; produto = cinza médio.
% FREEZE respeitado: todos os números são os da própria Seção 3.2.3
% (>=5, >=2, 0,13%, 621, 620, 715, 714, 231.490), nenhum novo.
% ============================================================================
\begin{tikzpicture}[
    x=1cm, y=1cm,
    caixa/.style={draw, rounded corners=2pt, align=center, text width=42mm,
                  minimum height=11mm, inner sep=3pt, font=\footnotesize},
    entrada/.style={caixa, fill=gray!8},
    produto/.style={caixa, fill=gray!14},
    seta/.style={-{Stealth[length=2.4mm]}, thick},
    rot/.style={font=\scriptsize, align=center, inner sep=2pt}
]
    \node[entrada, text width=40mm] (base) at (4.6,0)
        {base corrigida pela auditoria (Seção~\ref{sec:metodo-dados-auditoria}):
         linhas de descrição curta $+$ rótulo};
    \node[caixa, text width=46mm] (norm) at (4.6,-2.4)
        {normalização mínima: minúsculas $+$ remoção de acentos (sem efeito
         nas contagens desta base); rótulos idem; variações lexicais ficam
         com o tokenizador de subpalavras};
    \draw[seta] (base) -- (norm);

    % ------- bifurcação: a mesma base, contada de dois jeitos -------------
    % (o título de cada coluna é a 1ª linha da própria caixa: títulos soltos
    % eram atravessados pelas setas da bifurcação — achado visual)
    \node[caixa, text width=44mm] (filtro5) at (1.4,-5.3)
        {\textbf{espaço dos oráculos} (sobre as linhas): classes com $<5$
         instâncias agrupadas na sentinela \texttt{\_rare\_} ($0{,}13\%$ das
         linhas)};
    \node[caixa, text width=44mm] (dedup) at (8.2,-5.3)
        {\textbf{visão populacional}:\\deduplicação (chave: descrição
         aparada, em minúsculas; mantém a 1ª ocorrência) $+$ filtro brando
         $\ge 2$: 715 classes em linhas, 714 permanecem};

    \node[produto, text width=44mm] (catsch) at (1.4,-8.3)
        {\textit{CategorySchema}: espaço fechado\\
         com \textbf{621 categorias}\\
         (620 mais a sentinela \texttt{\_rare\_});\\
         fonte única da restrição de saída\\
         dos oráculos (Seção~\ref{sec:metodo-oraculo-instrumentacao})};
    \node[produto, text width=44mm] (pop) at (8.2,-8.3)
        {\textbf{231.490 textos únicos} em \textbf{714 classes presentes},
         sem \texttt{\_rare\_} (cauda longa preservada): experimentos do
         Capítulo~\ref{ch:resultados-falco}};

    % setas da bifurcação desenhadas depois dos nós (referência-adiante)
    \draw[seta] (norm.south west) -- (filtro5.north);
    \draw[seta] (norm.south east) -- (dedup.north);
    \draw[seta] (filtro5) -- (catsch);
    \draw[seta] (dedup) -- (pop);

    \node[rot, anchor=west, text width=118mm, align=left] at (-0.9,-10.4)
        {Duas chaves de texto convivem: o censo da auditoria compara
         descrições com os espaços internos colapsados; a deduplicação usa a
         descrição apenas aparada e em minúsculas. A receita executável
         \texttt{scripts/check\_dataset.py} reproduz todas as contagens.};
\end{tikzpicture}
\end{document}
%   (no modo standalone, troque as \ref{...} por texto fixo ou defina-as)
% Requer apenas arrows.meta e positioning (já em packages.tex). Legível em
% P&B: entrada = cinza claro; operação = branco; produto = cinza médio.
% FREEZE respeitado: todos os números são os da própria Seção 3.2.3
% (>=5, >=2, 0,13%, 621, 620, 715, 714, 231.490), nenhum novo.
% ============================================================================
\begin{tikzpicture}[
    x=1cm, y=1cm,
    caixa/.style={draw, rounded corners=2pt, align=center, text width=42mm,
                  minimum height=11mm, inner sep=3pt, font=\footnotesize},
    entrada/.style={caixa, fill=gray!8},
    produto/.style={caixa, fill=gray!14},
    seta/.style={-{Stealth[length=2.4mm]}, thick},
    rot/.style={font=\scriptsize, align=center, inner sep=2pt}
]
    \node[entrada, text width=40mm] (base) at (4.6,0)
        {base corrigida pela auditoria (Seção~\ref{sec:metodo-dados-auditoria}):
         linhas de descrição curta $+$ rótulo};
    \node[caixa, text width=46mm] (norm) at (4.6,-2.4)
        {normalização mínima: minúsculas $+$ remoção de acentos (sem efeito
         nas contagens desta base); rótulos idem; variações lexicais ficam
         com o tokenizador de subpalavras};
    \draw[seta] (base) -- (norm);

    % ------- bifurcação: a mesma base, contada de dois jeitos -------------
    % (o título de cada coluna é a 1ª linha da própria caixa: títulos soltos
    % eram atravessados pelas setas da bifurcação — achado visual)
    \node[caixa, text width=44mm] (filtro5) at (1.4,-5.3)
        {\textbf{espaço dos oráculos} (sobre as linhas): classes com $<5$
         instâncias agrupadas na sentinela \texttt{\_rare\_} ($0{,}13\%$ das
         linhas)};
    \node[caixa, text width=44mm] (dedup) at (8.2,-5.3)
        {\textbf{visão populacional}:\\deduplicação (chave: descrição
         aparada, em minúsculas; mantém a 1ª ocorrência) $+$ filtro brando
         $\ge 2$: 715 classes em linhas, 714 permanecem};

    \node[produto, text width=44mm] (catsch) at (1.4,-8.3)
        {\textit{CategorySchema}: espaço fechado\\
         com \textbf{621 categorias}\\
         (620 mais a sentinela \texttt{\_rare\_});\\
         fonte única da restrição de saída\\
         dos oráculos (Seção~\ref{sec:metodo-oraculo-instrumentacao})};
    \node[produto, text width=44mm] (pop) at (8.2,-8.3)
        {\textbf{231.490 textos únicos} em \textbf{714 classes presentes},
         sem \texttt{\_rare\_} (cauda longa preservada): experimentos do
         Capítulo~\ref{ch:resultados-falco}};

    % setas da bifurcação desenhadas depois dos nós (referência-adiante)
    \draw[seta] (norm.south west) -- (filtro5.north);
    \draw[seta] (norm.south east) -- (dedup.north);
    \draw[seta] (filtro5) -- (catsch);
    \draw[seta] (dedup) -- (pop);

    \node[rot, anchor=west, text width=118mm, align=left] at (-0.9,-10.4)
        {Duas chaves de texto convivem: o censo da auditoria compara
         descrições com os espaços internos colapsados; a deduplicação usa a
         descrição apenas aparada e em minúsculas. A receita executável
         \texttt{scripts/check\_dataset.py} reproduz todas as contagens.};
\end{tikzpicture}
\end{document}
%   (no modo standalone, troque as \ref{...} por texto fixo ou defina-as)
% Requer apenas arrows.meta e positioning (já em packages.tex). Legível em
% P&B: entrada = cinza claro; operação = branco; produto = cinza médio.
% FREEZE respeitado: todos os números são os da própria Seção 3.2.3
% (>=5, >=2, 0,13%, 621, 620, 715, 714, 231.490), nenhum novo.
% ============================================================================
\begin{tikzpicture}[
    x=1cm, y=1cm,
    caixa/.style={draw, rounded corners=2pt, align=center, text width=42mm,
                  minimum height=11mm, inner sep=3pt, font=\footnotesize},
    entrada/.style={caixa, fill=gray!8},
    produto/.style={caixa, fill=gray!14},
    seta/.style={-{Stealth[length=2.4mm]}, thick},
    rot/.style={font=\scriptsize, align=center, inner sep=2pt}
]
    \node[entrada, text width=40mm] (base) at (4.6,0)
        {base corrigida pela auditoria (Seção~\ref{sec:metodo-dados-auditoria}):
         linhas de descrição curta $+$ rótulo};
    \node[caixa, text width=46mm] (norm) at (4.6,-2.4)
        {normalização mínima: minúsculas $+$ remoção de acentos (sem efeito
         nas contagens desta base); rótulos idem; variações lexicais ficam
         com o tokenizador de subpalavras};
    \draw[seta] (base) -- (norm);

    % ------- bifurcação: a mesma base, contada de dois jeitos -------------
    % (o título de cada coluna é a 1ª linha da própria caixa: títulos soltos
    % eram atravessados pelas setas da bifurcação — achado visual)
    \node[caixa, text width=44mm] (filtro5) at (1.4,-5.3)
        {\textbf{espaço dos oráculos} (sobre as linhas): classes com $<5$
         instâncias agrupadas na sentinela \texttt{\_rare\_} ($0{,}13\%$ das
         linhas)};
    \node[caixa, text width=44mm] (dedup) at (8.2,-5.3)
        {\textbf{visão populacional}:\\deduplicação (chave: descrição
         aparada, em minúsculas; mantém a 1ª ocorrência) $+$ filtro brando
         $\ge 2$: 715 classes em linhas, 714 permanecem};

    \node[produto, text width=44mm] (catsch) at (1.4,-8.3)
        {\textit{CategorySchema}: espaço fechado\\
         com \textbf{621 categorias}\\
         (620 mais a sentinela \texttt{\_rare\_});\\
         fonte única da restrição de saída\\
         dos oráculos (Seção~\ref{sec:metodo-oraculo-instrumentacao})};
    \node[produto, text width=44mm] (pop) at (8.2,-8.3)
        {\textbf{231.490 textos únicos} em \textbf{714 classes presentes},
         sem \texttt{\_rare\_} (cauda longa preservada): experimentos do
         Capítulo~\ref{ch:resultados-falco}};

    % setas da bifurcação desenhadas depois dos nós (referência-adiante)
    \draw[seta] (norm.south west) -- (filtro5.north);
    \draw[seta] (norm.south east) -- (dedup.north);
    \draw[seta] (filtro5) -- (catsch);
    \draw[seta] (dedup) -- (pop);

    \node[rot, anchor=west, text width=118mm, align=left] at (-0.9,-10.4)
        {Duas chaves de texto convivem: o censo da auditoria compara
         descrições com os espaços internos colapsados; a deduplicação usa a
         descrição apenas aparada e em minúsculas. A receita executável
         \texttt{scripts/check\_dataset.py} reproduz todas as contagens.};
\end{tikzpicture}
\end{document}
%   (no modo standalone, troque as \ref{...} por texto fixo ou defina-as)
% Requer apenas arrows.meta e positioning (já em packages.tex). Legível em
% P&B: entrada = cinza claro; operação = branco; produto = cinza médio.
% FREEZE respeitado: todos os números são os da própria Seção 3.2.3
% (>=5, >=2, 0,13%, 621, 620, 715, 714, 231.490), nenhum novo.
% ============================================================================
\begin{tikzpicture}[
    x=1cm, y=1cm,
    caixa/.style={draw, rounded corners=2pt, align=center, text width=42mm,
                  minimum height=11mm, inner sep=3pt, font=\footnotesize},
    entrada/.style={caixa, fill=gray!8},
    produto/.style={caixa, fill=gray!14},
    seta/.style={-{Stealth[length=2.4mm]}, thick},
    rot/.style={font=\scriptsize, align=center, inner sep=2pt}
]
    \node[entrada, text width=40mm] (base) at (4.6,0)
        {base corrigida pela auditoria (Seção~\ref{sec:metodo-dados-auditoria}):
         linhas de descrição curta $+$ rótulo};
    \node[caixa, text width=46mm] (norm) at (4.6,-2.4)
        {normalização mínima: minúsculas $+$ remoção de acentos (sem efeito
         nas contagens desta base); rótulos idem; variações lexicais ficam
         com o tokenizador de subpalavras};
    \draw[seta] (base) -- (norm);

    % ------- bifurcação: a mesma base, contada de dois jeitos -------------
    % (o título de cada coluna é a 1ª linha da própria caixa: títulos soltos
    % eram atravessados pelas setas da bifurcação — achado visual)
    \node[caixa, text width=44mm] (filtro5) at (1.4,-5.3)
        {\textbf{espaço dos oráculos} (sobre as linhas): classes com $<5$
         instâncias agrupadas na sentinela \texttt{\_rare\_} ($0{,}13\%$ das
         linhas)};
    \node[caixa, text width=44mm] (dedup) at (8.2,-5.3)
        {\textbf{visão populacional}:\\deduplicação (chave: descrição
         aparada, em minúsculas; mantém a 1ª ocorrência) $+$ filtro brando
         $\ge 2$: 715 classes em linhas, 714 permanecem};

    \node[produto, text width=44mm] (catsch) at (1.4,-8.3)
        {\textit{CategorySchema}: espaço fechado\\
         com \textbf{621 categorias}\\
         (620 mais a sentinela \texttt{\_rare\_});\\
         fonte única da restrição de saída\\
         dos oráculos (Seção~\ref{sec:metodo-oraculo-instrumentacao})};
    \node[produto, text width=44mm] (pop) at (8.2,-8.3)
        {\textbf{231.490 textos únicos} em \textbf{714 classes presentes},
         sem \texttt{\_rare\_} (cauda longa preservada): experimentos do
         Capítulo~\ref{ch:resultados-falco}};

    % setas da bifurcação desenhadas depois dos nós (referência-adiante)
    \draw[seta] (norm.south west) -- (filtro5.north);
    \draw[seta] (norm.south east) -- (dedup.north);
    \draw[seta] (filtro5) -- (catsch);
    \draw[seta] (dedup) -- (pop);

    \node[rot, anchor=west, text width=118mm, align=left] at (-0.9,-10.4)
        {Duas chaves de texto convivem: o censo da auditoria compara
         descrições com os espaços internos colapsados; a deduplicação usa a
         descrição apenas aparada e em minúsculas. A receita executável
         \texttt{scripts/check\_dataset.py} reproduz todas as contagens.};
\end{tikzpicture}
